\subsection{Uniform Continuity and Compactness}

\begin{definition} \label{def:uniform-continuity} \hlt{Uniform Continuity}\\
Let $(X,d_X)$ and $(Y,d_Y)$ be metric spaces, and let $f:X\to Y$.\\
The map $f$ is \hlt{uniformly continuous} iff $(\forall\,\varepsilon>0)(\exists\,\delta>0)(\forall\,p,q\in X)\,[d_X(p,q)<\delta\Longrightarrow d_Y(f(p),f(q))<\varepsilon]$.\\
\hlt{Key Insight}: The choice of $\delta$ may depend on $\varepsilon$, but not on either point in the pair.
\end{definition}

\begin{principle} \label{princ:continuity-challenge-response} \hlt{Challenge--Response Order}\\
Limits and continuity depend on quantifier order. Definition \ref{def:function-limit-metric} answers a range challenge $\varepsilon$ by a domain radius $\delta$ at the base point.\\
Definition \ref{def:uniform-continuity} strengthens this pattern by requiring $\delta$ to be independent of the domain point.
\end{principle}

\begin{proposition} \label{prop:uniform-continuity-sequential-test} \hlt{Sequential Test for Uniform Continuity}\\
Let $X,Y$ be metric spaces, $f:X\to Y$. Then $f$ uniformly continuous iff $d_Y(f(p_n),f(q_n))\to0$ when $(p_n),(q_n)\subseteq X$, $d_X(p_n,q_n)\to0$.
\begin{proof}
\begin{enumerate}[label=\roman*.]
\setlength{\itemsep}{0pt}
\item[$(\Rightarrow)$] Let $\varepsilon>0$. By uniform continuity, choose $\delta>0$ such that $x,y\in X$ and $d_X(x,y)<\delta$ imply $d_Y(f(x),f(y))<\varepsilon$.\\
Since $d_X(p_n,q_n)\to0$, eventually $d_X(p_n,q_n)<\delta$, and hence eventually $d_Y(f(p_n),f(q_n))<\varepsilon$.
\item[$(\Leftarrow)$] Prove contrapositive. If $f$ not uniformly continuous, $(\exists\varepsilon_0>0)(\forall n\geq1)(\exists p_n,q_n\in X)\,[d_X(p_n,q_n)<\frac{1}{n}\land d_Y(f(p_n),f(q_n))\geq\varepsilon_0]$.\\
$\frac{1}{n}\to0$ by Theorem \ref{thm:special-sequence-limits}(i), $d_X(p_n,q_n)\to0$ by Theorem \ref{thm:pinching-theorem-sequences}; $d_Y(f(p_n),f(q_n))\not\to0$, contradicting sequential condition.
\end{enumerate}
\hlt{Structural Commentary}: The forward estimate transfers to pairs; the negation selects a separated pair at each scale.\\
\hlt{Key Insight}: Failed uniformity produces close inputs whose outputs stay apart.
\end{proof}
\end{proposition}

\begin{definition} \label{def:lipschitz-map} \hlt{Lipschitz Map}\\
Let $(X,d_X)$ and $(Y,d_Y)$ be metric spaces, $f:X\to Y$, $L\geq0$. The map $f$ is \hlt{$L$-Lipschitz} if $(\forall\,p,q\in X)\ d_Y(f(p),f(q))\leq Ld_X(p,q)$.\\
It is \hlt{Lipschitz} if $(\exists\,L\geq0)$ such that it is $L$-Lipschitz.\\
\hlt{Key Insight}: Lipschitz control is one global linear modulus of continuity.
\end{definition}

\begin{proposition} \label{prop:lipschitz-implies-uniform-continuity} \hlt{Lipschitz Implies Uniform Continuity}\\
Let $(X,d_X)$ and $(Y,d_Y)$ be metric spaces, and let $f:X\to Y$. If $f$ is Lipschitz, then $f$ is uniformly continuous and continuous.
\begin{proof}
Choose $L\geq0$ such that $f$ is $L$-Lipschitz by Definition \ref{def:lipschitz-map}, and let $\varepsilon>0$.\\
If $L=0$, set $\delta\defeq1$; then $(\forall\,p,q\in X)\ d_Y(f(p),f(q))=0<\varepsilon$.\\
If $L>0$, take $\delta\defeq\frac{\varepsilon}{L}$; if $p,q\in X$ and $d_X(p,q)<\delta$, then $d_Y(f(p),f(q))\leq Ld_X(p,q)<\varepsilon$. Thus $f$ is uniformly continuous.\\
Fix $a\in X$. If $p\in X$ and $d_X(p,a)<\delta$, the uniform estimate with $q=a$ verifies continuity at $a$.\\
\hlt{Structural Commentary}: The cases $L=0$ and $L>0$ provide a point-independent $\delta$; specialization proves continuity.\\
\hlt{Key Insight}: The linear bound supplies a global modulus.
\end{proof}
\end{proposition}

\begin{theorem} \label{thm:compact-continuity-uniform-continuity} \hlt{Continuity on Compact Sets Is Uniform}\\
Let $(X,d_X)$, $(Y,d_Y)$ be metric spaces. If $f:X\to Y$ is continuous and $X$ is compact, then $f$ is uniformly continuous.
\begin{proof}
If $X=\emptyset$, the conclusion is vacuous. Assume $X\neq\emptyset$, and fix $\varepsilon>0$. For $a\in X$, set $U_a\defeq f^{-1}(B_{\frac{\varepsilon}{2}}(f(a)))$.\\
Each $U_a$ is open by Theorems \ref{thm:every-neighborhood-open} and \ref{thm:continuity-open-preimages}, while $a\in U_a$. Hence $\mathcal{U}\defeq\{U_a:a\in X\}$ is an open cover of $X$.\\
Choose Lebesgue number $\delta>0$ for $\mathcal{U}$ by Lemma \ref{lem:lebesgue-covering-lemma}(ii). If $p,q\in X$, $d_X(p,q)<\delta$, choose $a\in X$, $B_\delta(p)\subseteq U_a$ by Definition \ref{def:lebesgue-number-open-cover}.\\
Then $p,q\in U_a$, so $d_Y(f(p),f(a))<\frac{\varepsilon}{2}$ and $d_Y(f(q),f(a))<\frac{\varepsilon}{2}$. Thus $d_Y(f(p),f(q))<\varepsilon$ by the triangle inequality.\\
\hlt{Structural Commentary}: Output balls induce an inverse-image cover; its Lebesgue number controls every input pair.\\
\hlt{Key Insight}: One cover-scale replaces all point-dependent continuity radii.
\end{proof}
\end{theorem}

\begin{counterexample} \label{cex:continuity-uniform-lipschitz-reversals} \hlt{The Reverse Implications Fail}
\begin{enumerate}[label=(\roman*)]
\setlength{\itemsep}{0pt}
\item The identity on $\R$ is an isometry, so it is continuous by Theorem \ref{thm:metric-isometry-basic-properties}(ii).\\
Thus $f:\R\to\R$, $f(x)\defeq x^2$, is continuous by Theorem \ref{thm:continuous-algebra-and-components}(i).\\
For $p_n\defeq n$ and $q_n\defeq n+\frac{1}{n}$, $\abs{p_n-q_n}=\frac{1}{n}\to0$, whereas $\abs{f(q_n)-f(p_n)}=2+\frac{1}{n^2}\nrightarrow0$.\\
Thus $f$ is not uniformly continuous by Proposition \ref{prop:uniform-continuity-sequential-test}.
\item Define $g:[0,1]\to\R$ by $g(x)\defeq\sqrt{x}$. For $\varepsilon>0$, the inequality $\abs{\sqrt{x}-\sqrt{y}}\leq\sqrt{\abs{x-y}}$ shows that $\delta\defeq\varepsilon^2$ verifies continuity.\\
The set $[0,1]$ is compact by Theorem \ref{thm:heine-borel-rk}. Hence $g$ is uniformly continuous by Theorem \ref{thm:compact-continuity-uniform-continuity}.\\
If $g$ were $L$-Lipschitz, then $(\forall n\in\N)\ \frac{1}{n}=\abs{g(\frac{1}{n^2})-g(0)}\leq\frac{L}{n^2}$, so $n\leq L$.\\
Choose $n\in\N$ with $n>L$ by Theorem \ref{thm:archimedean-and-rational-density}(i), a contradiction. Hence $g$ is not Lipschitz.
\end{enumerate}
\hlt{Structural Commentary}: The quadratic breaks uniformity; the square root has no linear modulus.\\
\hlt{Key Insight}: Compactness supplies a modulus, but not necessarily a linear one.
\end{counterexample}

\begin{theorem} \label{thm:compactness-hypotheses-are-essential} \hlt{Compactness Hypotheses Are Essential}\\
Let $E\subseteq\R$ be noncompact in the subspace metric. Then:
\begin{enumerate}[label=(\roman*)]
\setlength{\itemsep}{0pt}
\item $\exists f:E\to\R$ continuous and unbounded;
\item $\exists g:E\to\R$ continuous and bounded with no maximum;
\item if $E$ is bounded, then $\exists h:E\to\R$ continuous and not uniformly continuous.
\end{enumerate}
\begin{proof}
The inclusion $\iota:E\to\R$ is an isometry, so it is continuous by Theorem \ref{thm:metric-isometry-basic-properties}(ii).\\
For $c\in\R$, define $\kappa_c:E\to\R$ by $\kappa_c(x)\defeq c$. Since $(\forall p,x\in E)\ \abs{\kappa_c(x)-\kappa_c(p)}=0$, $\kappa_c$ is continuous by Definition \ref{def:continuity-at-a-point-and-on-a-set}.\\
Their denominators are nonzero on $E$, so the rational expressions below are continuous by Theorem \ref{thm:continuous-algebra-and-components}(i).
\begin{enumerate}[label=(\roman*)]
\setlength{\itemsep}{0pt}
\item If $E$ is unbounded, $f\defeq\iota$ is continuous and unbounded. Suppose $E$ is bounded.\\
Then $E$ not closed by Theorem \ref{thm:heine-borel-rk}; choose $x_0\in\overline{E}\setminus E$. Point $x_0$ is limit point of $E$ by Definition \ref{def:metric-closure}.\\
Define $f(x)\defeq\frac{1}{x-x_0}$. For $M>0$, choose $x\in E\cap B_{\frac{1}{M}}(x_0)$ by limit-point property of $x_0$. Then $\abs{f(x)}>M$, so $f$ is unbounded.
\item Suppose first that $E$ is unbounded, and define $g(x)\defeq\frac{x^2}{1+x^2}$. Then $(\forall x\in E)\ 0\leq g(x)<1$.\\
Given $\varepsilon>0$, choose $n\in\N$ with $n>\varepsilon^{-\frac{1}{2}}$ by Theorem \ref{thm:archimedean-and-rational-density}(i). Choose $x\in E$ with $\abs{x}>n$ by unboundedness.\\
Then $1-g(x)=\frac{1}{1+x^2}<\varepsilon$. Since $g<1$ on $E$, $\sup\limits_E g=1$, which is not attained. Suppose next that $E$ is bounded.\\
Choose $x_0\in\overline{E}\setminus E$ as in (i), define $g(x)\defeq\frac{1}{1+(x-x_0)^2}$. Then $(\forall x\in E)\ 0<g(x)<1$.\\
Given $\varepsilon>0$, $x\in E\cap B_{\sqrt{\varepsilon}}(x_0)$ by limit-point property of $x_0$; $1-g(x)=\frac{(x-x_0)^2}{1+(x-x_0)^2}<\varepsilon$. Hence $\sup\limits_E g=1$, which is not attained.
\item Suppose $E$ is bounded. Choose $x_0\in\overline{E}\setminus E$ as in (i), and define $h(x)\defeq\frac{1}{x-x_0}$.\\
For $\delta>0$, choose $x\in E\cap B_{\frac{\delta}{2}}(x_0)$ by limit-point property of $x_0$. Choose $y\in E$ with $\abs{y-x_0}<\min\{\frac{\delta}{2},\frac{1}{\abs{h(x)}+1}\}$ by same property.\\
Then $\abs{x-y}<\delta$ and $\abs{h(y)}>\abs{h(x)}+1$, so $\abs{h(y)-h(x)}>1$ by the reverse triangle inequality.\\
Thus $h$ is not uniformly continuous by Definition \ref{def:uniform-continuity}.
\end{enumerate}
\hlt{Structural Commentary}: Unbounded sets support escape functions; bounded nonclosed sets support singular boundary functions.\\
\hlt{Key Insight}: Every noncompact subset of $\R$ has an escape or boundary obstruction.
\end{proof}
\end{theorem}

\noindent Figure \ref{fig:uniform-continuity-delta-shrinks-one-over-x-squared} illustrates the boundary obstruction.\\
Near $0$, a fixed output tolerance for $x\mapsto\frac{1}{x^2}$ requires progressively smaller input radii.

\begin{visualizationfigure}[H]
  \centering
  \begin{tikzpicture}[
  scale=0.95,
  >=Stealth,
  every node/.style={font=\scriptsize},
  axis/.style={->, gray!70, line width=0.45pt},
  curve/.style={blue!70!black, line width=0.9pt},
  guide/.style={gray!55, dashed, line width=0.4pt},
  epsmark/.style={orange!80!black, line width=0.65pt},
  bracket/.style={teal!75!black, line width=0.8pt},
  sample/.style={blue!70!black}
]
  \draw[axis] (-0.25,0) -- (4.00,0) node[right] {$x$};
  \draw[axis] (0,-0.45) -- (0,2.65) node[above] {$y$};

  \draw[curve, domain=0.66:3.45, samples=120] plot ({\x},{1/(\x*\x)});

  \foreach \x/\name/\d/\dlabel in {2.85/$x_1$/0.46/$\delta_1$,1.75/$x_2$/0.26/$\delta_2$,1.03/$x_3$/0.11/$\delta_3$} {
    \pgfmathsetmacro{\y}{1/(\x*\x)}
    \fill[sample] (\x,\y) circle (0.035);
    \draw[guide] (\x,0) -- (\x,\y);
    \node[above] at (\x,0.03) {\name};
    \draw[bracket] ({\x-\d},-0.12) -- ({\x+\d},-0.12);
    \draw[bracket] ({\x-\d},-0.18) -- ({\x-\d},-0.06);
    \draw[bracket] ({\x+\d},-0.18) -- ({\x+\d},-0.06);
    \node[teal!75!black, below] at (\x,-0.19) {\dlabel};
  }

  \draw[guide] (3.35,0.45) -- (3.70,0.45);
  \draw[guide] (3.35,1.20) -- (3.70,1.20);
  \draw[<->, epsmark] (3.60,0.45) -- (3.60,1.20);
  \node[orange!80!black, right] at (3.60,0.825) {$2\varepsilon$};
\end{tikzpicture}

  \caption{Dimensionless $y=\frac{1}{x^2}$ graph for $x>0$: input radii $\delta_i$ shrink as $x_i\to0$; orange marks the fixed $2\varepsilon$ output tolerance.}
  \label{fig:uniform-continuity-delta-shrinks-one-over-x-squared}
\end{visualizationfigure}

\begin{proofmethod} \label{meth:continuity-proof-patterns} \hlt{Continuity Proof Patterns}\\
Let $(X,d_X)$ and $(Y,d_Y)$ be metric spaces, $E\subseteq X$, $f:E\to Y$, $p\in X$ and $q\in Y$.
\begin{enumerate}[label=(\roman*)]
\setlength{\itemsep}{0pt}
\item \hlt{$\varepsilon$--$\delta$ at $p$}: for $p\in E$, prove $(\forall\,\varepsilon>0)(\exists\,\delta>0)(\forall\,x\in E)\,[d_X(x,p)<\delta\Longrightarrow d_Y(f(x),f(p))<\varepsilon]$.
\item \hlt{Limit at $p$}: for $p\in E^{\operatorname{acc}}$, replace antecedent by $0<d_X(x,p)<\delta$ and target by $q$.\\
Equivalently, test sequences with $x_n\to p$ and $x_n\neq p$.
\item \hlt{Topological continuity}: prove that open subsets of $Y$ have relatively open preimages in $E$.\\
Equivalently, prove that closed subsets of $Y$ have relatively closed preimages.
\item \hlt{Compact-image argument}: for compact $E$, pull an open cover $\{V_\alpha\}$ of $f(E)$ back to $\{f^{-1}(V_\alpha)\}$ of $E$.\\
Extract a finite subcover of $E$ and push it forward to $f(E)$.
\item \hlt{Uniform continuity}: for compact $E$ and $\varepsilon>0$, form $U_a\defeq f^{-1}(B_{\frac{\varepsilon}{2}}(f(a)))$ for $a\in E$.\\
Choose a Lebesgue number $\delta>0$ by Lemma \ref{lem:lebesgue-covering-lemma}(ii).
\end{enumerate}
\hlt{Key Insight}: Match the target property to pointwise, sequential, inverse-image, or compact-cover control.
\end{proofmethod}
